\begin{figure}[htbp]
\centering
\begin{tikzpicture}[font=\small,
 box/.style={draw=black!55,rounded corners,inner sep=9pt,align=center},
 >=Stealth]
 \node (parent) at (0,1.4) {an admissible parent \(\mathbf I\)};
 \node[box,fill=blue!5] (plus) at (-2.5,0)
 {\(s_0=+1\)\\[3pt]
 \((+,+,+)\qquad(+,-,-)\)};
 \node[box,fill=blue!5] (minus) at (2.5,0)
 {\(s_0=-1\)\\[3pt]
 \((-,+,-)\qquad(-,-,+)\)};
 \draw[->] (parent) -- (plus);
 \draw[->] (parent) -- (minus);
 \node[align=center] at (0,-1.25)
 {\(\E[s_1\mid s_0]=\E[s_2\mid s_0]=0\),\qquad
   \(s_1s_2=s_0\).};
\end{tikzpicture}
\caption{The four admissible child sign triples, grouped by the
first sign. Conditional on \(s_0\), the two possibilities in its box
have equal probability. Each neighboring sign is unbiased, but their
product is fixed. The sum of the two incident-edge Grams in
\eqref{eq:child-gram} uses only these marginal averages.}
\label{fig:children}
\end{figure}
